\documentclass{article}
\usepackage[UTF8]{ctex}
\usepackage{amsmath,amssymb,amsfonts}
\usepackage{enumitem}
\usepackage{geometry}
\geometry{a4paper, margin=1in}
% === 错题本排版优化（不参与题目标识） ===
\setlength{\parindent}{0pt}
\setlength{\parskip}{0.35em plus 0.10em minus 0.08em}
\setlength{\emergencystretch}{3em}
\setlist{itemsep=0.25em, topsep=0.40em, parsep=0pt}
\allowdisplaybreaks[2]
% === 错题本排版优化结束 ===
\begin{document}
\section*{复旦大学2025年高等代数真题}
\section*{31．复旦大学2025年高等代数部分（回忆版）}
一．填空题．每题5分，共25分．\\

\subsection{题目1}
1．已知 \(A=\left(\begin{array}{cccc}2 & 1 & \cdots & 1 \\ 1 & 2 & \cdots & 1 \\ \vdots & \vdots & & \vdots \\ 1 & 1 & \cdots & 2\end{array}\right)\) ，求 \(A^{2024}=\) \(\_\_\_\_\)．\\
\(\_\_\_\_\)

\subsection{题目2}
2．\(K\) 为数域，定义 \(f: K^{n \times n} \rightarrow K\) 为线性映射，且 \(f\) 满足对相似的矩阵有相同的函数值，求所有可能的 \(f\) \(\_\_\_\_\)。

\subsection{题目3}
3．定义 \(f: K^{m \times n} \rightarrow K^{m \times n}\) 为线性映射，且 \(f(X)=X A\) ，其中 \(A\) 是一个秩为 \(r\) 的 \(n \times n\) 矩阵，求 \(f\)的秩为 \(\_\_\_\_\)．

\subsection{题目4}
4．已知矩阵 \(A=\left(\begin{array}{cc}2024 & 1 \\ 0 & 2024\end{array}\right)\) ，线性映射 \(\varphi: K^{2 \times 2} \rightarrow K^{2 \times 2}\) 满足 \(\varphi(X)=J \times J\) ，其中 \(X \in M_{2}(K)\)，求 \(\varphi\) 的对应矩阵的若尔当标准型 \(\_\_\_\_\)．

\subsection{题目5}
5．由1，2，3 组成的3阶行列式中值不是3的倍数的行列式有 \(\_\_\_\_\)．

二．（20 分）设 \(V, W\) 为数域上的非零向量空间，其中 \(V\) 是有限维的．\(\alpha_{1}, \alpha_{2}, \cdots, \alpha_{n} \in V\) 满足对任意的\(\beta_{1}, \beta_{2}, \cdots, \beta_{n} \in W\) ，存在唯一的线性映射 \(\varphi: V \rightarrow W\) ，使得

\[
\varphi\left(\alpha_{i}\right)=\beta_{i}, i=1,2, \cdots, n .
\]

证明：\(\alpha_{1}, \alpha_{2}, \cdots, \alpha_{n}\) 是 \(V\) 的基．

三。（20分）设 \(V\) 是由 \(\mathbb{C}^{n \times n}\) 上的所有置换矩阵张成的子空间，所谓置换矩阵，是指每行每列恰好有一个元素为 1 ，其余元素均为 0 的矩阵。求 \(\operatorname{dim} V\) 。

四．（20 分）\(f\) 是 \(V\) 上实对称双线性函数，正负惯性指数分别为 \(1, n-1\) 。

\begin{itemize}
  \item[（1）] 若 \(f(\alpha, \alpha)>0, f(\alpha, \beta)=0\) ，证明：\(f(\beta, \beta) \leq 0\) ，等号成立当且仅当 \(\beta=0\) 。
  \item[（2）] 若 \(f(\alpha, \alpha)=f(\alpha, \beta)=f(\beta, \beta)=0\) ，证明：\(\alpha, \beta\) 线性相关．
\end{itemize}

五。（20分）\(A\) 为 \(n\) 阶可逆矩阵，\(H\) 为 \(n\) 阶实对称正定矩阵。证明：存在 \(n\) 阶正交矩阵 \(P\) ，使得

\[
\left(H^{\mathrm{T}}+E\right) A^{\mathrm{T}} A(E+H)-P^{\mathrm{T}} A^{\mathrm{T}} A P
\]

为实对称矩阵，其中 \(E\) 为单位矩阵。
\end{document}
